\documentclass[11pt]{article}
\usepackage[margin=1in]{geometry}
\usepackage[T1]{fontenc}
\usepackage{lmodern}
\usepackage{amsmath,amssymb,amsthm}
\usepackage[hidelinks]{hyperref}
\usepackage{microtype}
\raggedright

\newtheorem*{claim}{Claim}

\begin{document}

\begin{center}
{\Large Literature-implied answer for arXiv:2412.20282}\\[4pt]
{\large Defective logarithmic Sobolev inequalities survive bounded
ground-state perturbations}
\end{center}

\paragraph{Status.}
\textbf{literature\_implied\_answer (full bounded-perturbation case).}
This is an agent-identified composition of published theorems, not a new proof
produced by the run.

\paragraph{Source question.}
In Remark~6.19, PDF pages 69--70, Gross asks whether, for a bounded potential
$V_1$, the ground-state measure of $-\Delta+V_0+V_1$ satisfies a defective
logarithmic Sobolev inequality whenever the ground-state measure of
$-\Delta+V_0$ does.  He concludes:
\begin{quote}
As to whether perturbation of a DLSI yields another DLSI under some
conditions on the perturbing potential is still an open question.
\end{quote}

Write the original ground-state probability measure as $\mu$ and its
Dirichlet form as $\mathcal E$.  Normalize the DLSI as
\begin{equation}
 \operatorname{Ent}_{\mu}(f^2)
 \leq 2C\,\mathcal E(f,f)+D\lVert f\rVert_{L^2(\mu)}^2.
 \label{eq:dlsi}
\end{equation}

\begin{claim}
Assume that the symmetric Markov semigroup associated with
$(\mathcal E,\mu)$ is ergodic.  If \eqref{eq:dlsi} holds and $W$ is a bounded real
potential, then the ground-state measure obtained by perturbing with $W$
satisfies a logarithmic Sobolev inequality with zero defect.  In particular,
it satisfies a DLSI.  The same conclusion holds for an unbounded $W$ whenever
Gross's exponential-integrability hypotheses hold with respect to the
tightened LSI constant.
\end{claim}

In the Schrodinger setting of the source, uniqueness of the original ground
state gives the required ergodicity after the ground-state transform.

\paragraph{Derivation.}
A DLSI is equivalent to hyperboundedness of the associated Markov semigroup.
Consequently, for some $T>0$ and $p>2$, the self-adjoint Markov operator
$P_T$ is bounded from $L^2(\mu)$ to $L^p(\mu)$.  Ergodicity of the semigroup
implies ergodicity of $P_T$ by spectral calculus.  Miclo's Theorem~1
\cite{Miclo} therefore gives a spectral gap.  Equivalently, there is a finite
constant $C_P$ such that
\[
 \operatorname{Var}_{\mu}(f)\leq C_P\mathcal E(f,f).
\]
Gross quotes the Rothaus tightening statement as Proposition~7.16, PDF
pages 85--86 \cite{Gross}: combining this Poincare inequality with
\eqref{eq:dlsi} gives
\begin{equation}
 \operatorname{Ent}_{\mu}(f^2)
 \leq 2c_0\mathcal E(f,f),
 \qquad c_0=C+C_P\left(\frac D2+1\right).
 \label{eq:lsi}
\end{equation}

Now apply Gross's main perturbation theorem, Theorem~2.2, PDF page 10, to
the zero-defect LSI \eqref{eq:lsi} and the potential $W$.  If $W$ is bounded, then for
every $\kappa>0$ and every $\nu>2c_0$,
\[
 \lVert e^W\rVert_{L^\kappa(\mu)}<\infty,
 \qquad
 \lVert e^{-W}\rVert_{L^\nu(\mu)}<\infty.
\]
The theorem produces a unique positive ground state for the perturbed
operator and a zero-defect LSI for its ground-state measure.  By the
consecutive ground-state transformation described in Gross's Section~8.1,
this is exactly the ground-state measure of $-\Delta+V_0+W$.  This proves the
claim.

The last step is unchanged for unbounded $W$ if the two displayed exponential
moments are finite for some $\kappa>0$ and $\nu>2c_0$.

\paragraph{Why this is a literature implication.}
Miclo proves the decisive hyperboundedness-to-gap theorem but predates
Gross's question and does not claim to answer it.  Gross himself cites Miclo
immediately after Remark~6.19 and separately records both Rothaus tightening
and the required perturbation theorem, but does not compose the three facts.
The relation to the stated question is therefore an identification, not an
independent new mathematical ingredient.

\paragraph{Limitations.}
Ergodicity is essential to Miclo's theorem, so the argument does not cover
arbitrary reducible forms.  The conversion from (1) to a gap is qualitative:
Miclo's Proposition~11 shows that no lower gap bound depending only on a
general hyperboundedness norm can exist.  Thus $c_0$ is explicit only when a
Poincare constant $C_P$ is independently available.

\paragraph{Search evidence.}
Cheap run indexes and bounded web/arXiv searches through 2026-08-09 used
arXiv:2412.20282, the exact question sentence, and combinations of
``defective logarithmic Sobolev'', ``bounded perturbation'', ``ground state'',
and ``Miclo''.  No later paper explicitly making this identification was
found.

\begin{thebibliography}{9}
\bibitem{Gross}
L. Gross,
\emph{Invariance of intrinsic hypercontractivity under perturbation of
Schrodinger operators},
Advanced Nonlinear Studies 25 (2025), no.~4, 951--1024;
arXiv:2412.20282v2.

\bibitem{Miclo}
L. Miclo,
\emph{On hyperboundedness and spectrum of Markov operators},
Inventiones Mathematicae 200 (2015), no.~1, 311--343,
doi:10.1007/s00222-014-0538-8.
\end{thebibliography}

\end{document}
